% =====================================================================
% Relatorio CDEG 2025/2026 - Esboco LaTeX
% Rendibilidade e risco de diferentes activos financeiros (2015-2025)
% Compilar com:  xelatex relatorio.tex   (duas vezes para o indice)
% =====================================================================
\documentclass[11pt,a4paper]{article}

% ---------- Pacotes -------------------------------------------------
\usepackage[T1]{fontenc}
\usepackage[utf8]{inputenc}
\usepackage[portuguese]{babel}
\usepackage[margin=2.5cm]{geometry}
\usepackage{graphicx}
\usepackage{booktabs}
\usepackage{array}
\usepackage{caption}
\usepackage{float}
\usepackage{amsmath,amssymb}
\usepackage{hyperref}
\usepackage{xcolor}
\usepackage{enumitem}
\usepackage{listings}
\usepackage{fancyhdr}

\hypersetup{colorlinks=true, linkcolor=black, citecolor=black,
            urlcolor=blue!60!black}

% ---------- Estilo do codigo R --------------------------------------
\definecolor{codebg}{RGB}{245,245,245}
\definecolor{codekw}{RGB}{0,80,160}
\definecolor{codecm}{RGB}{120,120,120}
\definecolor{codestr}{RGB}{160,40,40}

\lstdefinelanguage{Renhanced}{
  keywords={if,else,for,while,function,return,TRUE,FALSE,NULL,NA,
            library,require,source,c,list,data.frame,mean,median,sd,
            sum,length,seq,sapply,lapply,nrow,ncol,filter,group_by,
            summarise,mutate,arrange,select,lm,predict,sample,replicate,
            kmeans,scale,boot,quantile,cumprod,cummax,exp,log,rnorm},
  sensitive=true, comment=[l]{\#}, string=[b]",
}
\lstset{
  language=Renhanced,
  basicstyle=\ttfamily\footnotesize,
  backgroundcolor=\color{codebg},
  keywordstyle=\color{codekw}\bfseries,
  commentstyle=\color{codecm}\itshape,
  stringstyle=\color{codestr},
  frame=single, rulecolor=\color{black!20},
  breaklines=true, showstringspaces=false,
  numbers=left, numberstyle=\tiny\color{black!50},
  numbersep=8pt, xleftmargin=18pt, captionpos=b
}

% ---------- Cabecalho/rodape ---------------------------------------
\pagestyle{fancy}
\fancyhf{}
\fancyhead[L]{\small Ciencia dos Dados para Economia e Gestao}
\fancyhead[R]{\small FEUC -- 2025/2026}
\fancyfoot[C]{\thepage}
\renewcommand{\headrulewidth}{0.4pt}

% =====================================================================
\begin{document}

% ---------- CAPA ----------------------------------------------------
\begin{titlepage}
  \thispagestyle{empty}
  \centering
  \vspace*{2cm}

  {\Large\scshape Faculdade de Economia da Universidade de Coimbra\par}
  \vspace{0.4cm}
  {\large Ciencia dos Dados para Economia e Gestao\par}
  {\large 2025/2026\par}

  \vspace{3cm}
  \rule{\linewidth}{0.5pt}\\[0.6cm]
  {\Huge\bfseries Rendibilidade e risco\par}
  \vspace{0.2cm}
  {\Large de diferentes activos financeiros\\(2015--2025)\par}
  \rule{\linewidth}{0.5pt}\\[1cm]

  {\large\itshape Relatorio do projecto\par}

  \vfill
  {\large
    Nome 1 \quad -- \quad n.\textsuperscript{o} 0000000\\[0.2cm]
    Nome 2 \quad -- \quad n.\textsuperscript{o} 0000000\\[0.2cm]
    Nome 3 \quad -- \quad n.\textsuperscript{o} 0000000\\[0.2cm]
    Nome 4 \quad -- \quad n.\textsuperscript{o} 0000000\par
  }
  \vspace{1.5cm}
  {\large Docente: Pedro Bacao\par}
  {\large \today\par}
\end{titlepage}

\tableofcontents
\newpage

% =====================================================================
\section{Introducao}

Este relatorio acompanha a Shiny App desenvolvida no ambito da unidade
curricular \emph{Ciencia dos Dados para Economia e Gestao} (FEUC,
2025/2026). O tema escolhido foi a \textbf{comparacao da rendibilidade
e do risco de diferentes classes de activos financeiros} entre 2015 e
2025. Os objectivos sao: apresentar os dados; explicar a Shiny App;
documentar exemplos de codigo das tecnicas utilizadas (descritivas,
testes de hipoteses, bootstrap, Monte Carlo, regressao e
\emph{clustering}).

% =====================================================================
\section{Dados utilizados}

Os dados foram \textbf{descarregados do Yahoo Finance} (Close ajustado
diario) atraves do pacote Python \texttt{yfinance}, executado uma so
vez para gerar o ficheiro local \texttt{data/assets.csv}
(\textasciitilde{}35\,700 linhas, 6 colunas, 2.4\,MB). O universo de
activos esta na Tabela~\ref{tab:assets}.

\begin{table}[H]\centering
\caption{Universo de activos.}\label{tab:assets}
\begin{tabular}{lll}
\toprule
\textbf{Ticker} & \textbf{Nome} & \textbf{Classe} \\ \midrule
\verb|^GSPC|     & S\&P 500              & Accoes (EUA)         \\
\verb|^STOXX50E| & Euro Stoxx 50         & Accoes (Zona Euro)   \\
\verb|^FTSE|     & FTSE 100              & Accoes (Reino Unido) \\
\verb|^N225|     & Nikkei 225            & Accoes (Japao)       \\
\verb|PSI20.LS|  & PSI-20                & Accoes (Portugal)    \\
\verb|GC=F|      & Ouro (futuros)        & Materias-primas      \\
\verb|CL=F|      & Petroleo WTI          & Materias-primas      \\
\verb|BTC-USD|   & Bitcoin               & Criptomoedas         \\
\verb|ETH-USD|   & Ethereum              & Criptomoedas         \\
\verb|EURUSD=X|  & EUR / USD             & Cambio               \\
\verb|^TNX|      & Yield Treasuries 10Y  & Obrigacoes (yield)   \\
\verb|TLT|       & iShares 20+Y Treasury & Obrigacoes (preco)   \\
\bottomrule
\end{tabular}
\end{table}

% =====================================================================
\section{A Shiny App}

A app esta organizada em seis separadores (Tabela~\ref{tab:tabs}). A
barra lateral permite filtrar por activo, periodo e taxa sem risco.

\begin{table}[H]\centering
\caption{Separadores da Shiny App.}\label{tab:tabs}
\begin{tabular}{ll}
\toprule
\textbf{Separador} & \textbf{Conteudo} \\ \midrule
Dados             & Tabela interactiva, NA, estrutura \\
Exploracao        & Indice 100, rend. cumulativa, distribuicoes \\
Risco e retorno   & Descritivas, plano vol-retorno, correlacoes \\
Inferencia        & $t$-Welch, Jarque-Bera, bootstrap, Monte Carlo \\
Machine Learning  & $k$-means, dendrograma, regressao multipla \\
Sobre             & Descricao do projecto e do grupo \\
\bottomrule
\end{tabular}
\end{table}

\subsection{Como usar}
\begin{enumerate}[leftmargin=*]
  \item Abrir a pasta do projecto em RStudio.
  \item Executar \texttt{shiny::runApp()}.
  \item Escolher activos, periodo e taxa sem risco no painel lateral.
  \item Navegar pelos separadores; em \textbf{Inferencia} ajustar as
        replicas do bootstrap e o numero de trajectorias Monte Carlo;
        em \textbf{Machine Learning} variar $k$ e o activo alvo da
        regressao.
\end{enumerate}

% =====================================================================
\section{Exemplos de codigo}

\subsection{Calculo de rendibilidades e indice 100}
\begin{lstlisting}[caption={Indice base 100 e rendibilidades simples.}]
recompute_returns <- function(df) {
  df |> group_by(ticker) |>
    arrange(date, .by_group = TRUE) |>
    mutate(return = (close / lag(close) - 1) * 100) |> ungroup()
}

to_index_100 <- function(df) {
  df |> group_by(ticker) |> arrange(date, .by_group = TRUE) |>
    mutate(index100 = 100 * close / first(close)) |> ungroup()
}
\end{lstlisting}

\subsection{Anualizacao e racio de Sharpe}

Sob a hipotese de rendibilidades i.i.d., a media anualizada e
$\mu_a = \mu_d \cdot 252$ e a volatilidade anualizada
$\sigma_a = \sigma_d \sqrt{252}$. O Sharpe anualizado escreve-se
\[
  \mathrm{SR}_a = \frac{\mu_d - r_f^{(d)}}{\sigma_d}\sqrt{252},
  \qquad r_f^{(d)} = r_f^{(a)}/252.
\]

\subsection{Maximo drawdown}
\begin{lstlisting}[caption={Maximo drawdown a partir de uma serie de precos.}]
max_drawdown <- function(prices) {
  prices <- prices[!is.na(prices)]
  peak <- cummax(prices)
  min((prices - peak) / peak) * 100
}
\end{lstlisting}

\subsection{Bootstrap conjunto de media e Sharpe}
\begin{lstlisting}[caption={Bootstrap nao parametrico.}]
boot_mean_sharpe <- function(returns, R = 1000, freq = 252,
                             rf_annual_pct = 0) {
  rf_d <- rf_annual_pct / freq
  stat <- function(data, idx) {
    r <- data[idx]
    c(mean_ann = mean(r) * freq,
      sharpe   = (mean(r) - rf_d) / sd(r) * sqrt(freq))
  }
  boot::boot(returns, statistic = stat, R = R)
}
\end{lstlisting}

\subsection{Simulacao Monte Carlo (GBM)}

A trajectoria simulada de precos segue
$S_{t+1} = S_t \exp\!\big((\mu - \tfrac{1}{2}\sigma^2)\Delta t +
\sigma\sqrt{\Delta t}\,Z\big)$ com $Z \sim \mathcal{N}(0,1)$ e
$\Delta t = 1$ (um dia de negociacao).

\begin{lstlisting}[caption={GBM vectorizado.}]
mc_gbm <- function(S0, mu_d, sigma_d, n_steps = 252, n_paths = 200) {
  Z <- matrix(rnorm(n_steps * n_paths), nrow = n_steps)
  log_ret <- (mu_d - 0.5 * sigma_d^2) + sigma_d * Z
  rbind(rep(S0, n_paths), S0 * exp(apply(log_ret, 2, cumsum)))
}
\end{lstlisting}

\subsection{Clustering hierarquico sobre correlacoes}
\begin{lstlisting}[caption={Distancia 1 - |cor| e dendrograma.}]
cor_distance <- function(df) {
  W <- returns_wide(df) |> tidyr::drop_na()
  C <- cor(W[, -1], use = "complete.obs")
  as.dist(1 - abs(C))
}
hc <- hclust(cor_distance(assets), method = "average")
\end{lstlisting}

% --- Inserir aqui figuras exportadas do R em report/figs/ ----------
% \begin{figure}[H]\centering
%   \includegraphics[width=0.85\linewidth]{figs/index100.pdf}
%   \caption{Cotacoes normalizadas (base 100).}\label{fig:idx}
% \end{figure}

% =====================================================================
\section{Resultados e discussao}

\subsection{Risco e retorno}

Confirma-se a troca classica entre risco e retorno: as criptomoedas
exibem volatilidades anualizadas significativamente superiores as dos
indices accionistas, com Sharpe historico nao necessariamente
superior. O ouro e a EUR/USD apresentam volatilidade intermedia mas
correlacoes baixas com as accoes, caracteristicas tipicas de
diversificadores. As obrigacoes longas (\verb|TLT|) mostram
correlacao negativa com a yield (\verb|^TNX|), conforme esperado pela
duracao.

\subsection{Clustering e factores}

O \emph{clustering} hierarquico sobre $1 - |\rho|$ separa de forma
natural tres grupos de activos: (i) accoes desenvolvidas, (ii)
materias-primas e cambio, (iii) obrigacoes e cripto. A regressao
linear do PSI-20 sobre os restantes activos atribui a maior parte da
variancia ao Euro Stoxx 50, consistente com a integracao europeia.

% =====================================================================
\section{Limitacoes}
\begin{itemize}[leftmargin=*]
  \item Apenas dados diarios e Close ajustado; ignoramos
        \emph{intraday} e custos de transaccao.
  \item Periodo dominado por choques excepcionais (COVID-19, guerra
        na Ucrania, ciclo inflacionista 2022--2023).
  \item A simulacao GBM assume volatilidade constante e
        rendibilidades log-normais, hipoteses violadas pelas series
        reais (caudas pesadas, \emph{clustering} de volatilidade).
\end{itemize}

% =====================================================================
\section{Utilizacao de LLMs neste projecto}

\textbf{Modelos utilizados:} ChatGPT (GPT-4o e GPT-5) e GitHub Copilot
(com Claude Sonnet 4) ao longo do desenvolvimento.

\subsection*{Pedidos feitos aos LLMs}
A interaccao foi feita em ciclos curtos
(\emph{prompt $\rightarrow$ avaliacao critica $\rightarrow$ revisao}).
Os pedidos mais relevantes:
\begin{enumerate}[leftmargin=*]
  \item \emph{``Sugere uma estrutura de ficheiros para uma Shiny App
        em R com um relatorio R Markdown anexo, separando UI/server,
        funcoes estatisticas e funcoes de visualizacao.''}
  \item \emph{``Da-me um script Python que use \texttt{yfinance} para
        descarregar Close ajustado diario para varios tickers e
        gravar um CSV em formato longo (\texttt{date}, \texttt{ticker},
        \texttt{name}, \texttt{asset\_class}, \texttt{close},
        \texttt{return}).''}
  \item \emph{``Como anualizar correctamente a media e o
        desvio-padrao das rendibilidades diarias para construir um
        Sharpe ratio?''}
  \item \emph{``Implementa em R uma funcao para calcular o maximo
        drawdown a partir de um vector de precos sem dependencias
        externas.''}
  \item \emph{``Escreve uma funcao de bootstrap nao parametrico que
        devolva simultaneamente a media anualizada e o Sharpe, e
        mostra como extrair o IC percentil para cada uma.''}
  \item \emph{``Implementa em R uma simulacao GBM vectorizada que
        devolva uma matriz \texttt{n\_steps x n\_paths} de precos.''}
  \item \emph{``Sugere uma metrica de distancia entre activos baseada
        em correlacoes para usar com \texttt{hclust}, e explica
        porque \mbox{$1 - |\rho|$} faz sentido.''}
  \item \emph{``Reve este \texttt{ggplot2} do plano risco-retorno e
        melhora as anotacoes (cores por classe, \emph{labels} sem
        sobreposicao).''}
  \item \emph{``Como construir uma formula \texttt{lm()}
        programaticamente quando os nomes das colunas comecam por
        \texttt{\^{}} ou contem caracteres especiais (\texttt{=},
        \texttt{-})?''}
  \item \emph{``Faz uma revisao critica desta interpretacao da
        regressao do PSI-20 sobre os outros indices.''}
\end{enumerate}

\subsection*{Contribuicao dos LLMs para a versao final}
\begin{itemize}[leftmargin=*]
  \item Sugestao da divisao em \texttt{app.R} / \texttt{global.R} /
        \texttt{R/helpers.R} / \texttt{R/plots.R}.
  \item Estrutura inicial do \texttt{generate\_data.py} (loop sobre
        tickers, achatamento do MultiIndex devolvido pelo
        \texttt{yfinance}).
  \item Primeira versao das funcoes \texttt{boot\_mean\_sharpe()},
        \texttt{mc\_gbm()}, \texttt{max\_drawdown()},
        \texttt{cor\_distance()} e \texttt{factor\_lm()}, revistas
        para tratar \texttt{NA}, datas em falta e nomes de coluna
        especiais.
  \item Sugestoes de paleta e tema \texttt{ggplot2}.
  \item Esqueleto do cabecalho YAML do RMarkdown e da capa LaTeX.
\end{itemize}

\subsection*{Resultados nao incorporados (e porque)}
\begin{itemize}[leftmargin=*]
  \item \textbf{Migracao para \texttt{tidymodels}} --- descartada por
        aumentar significativamente o tempo de carregamento.
  \item \textbf{\texttt{shinydashboard}/\texttt{bslib}} --- preterido
        por \texttt{fluidPage} mais leve.
  \item \textbf{Modelos GARCH} para captar \emph{clustering} de
        volatilidade --- fora do ambito da unidade curricular.
  \item \textbf{Block bootstrap} para series temporais ---
        descartado por simplicidade pedagogica; ressalva-se no texto
        que o bootstrap i.i.d. subestima a incerteza.
  \item \textbf{Texto interpretativo automatico} sobre o crash de
        Marco de 2020 --- descartado por incluir afirmacoes
        factualmente imprecisas; reescrito a partir das estatisticas
        que efectivamente apuramos.
\end{itemize}

Toda a escolha do tema, a validacao dos dados, a interpretacao
financeira e a revisao final do texto e do codigo foram feitas pelos
membros do grupo.

% =====================================================================
\section*{Referencias}
\addcontentsline{toc}{section}{Referencias}
\begin{itemize}[leftmargin=*]
  \item Yahoo Finance --- \url{https://finance.yahoo.com}.
  \item Hyndman, R.J.; Athanasopoulos, G. (2021).
        \emph{Forecasting: Principles and Practice}. 3.\textsuperscript{a} ed.
        \url{https://otexts.com/fpp3/}
  \item James, G.; Witten, D.; Hastie, T.; Tibshirani, R. (2023).
        \emph{An Introduction to Statistical Learning with
        Applications in R}. 2.\textsuperscript{a} ed. Springer.
  \item Wickham, H. (2021). \emph{Mastering Shiny}. O'Reilly.
        \url{https://mastering-shiny.org}
  \item Wickham, H.; Grolemund, G. (2017). \emph{R for Data Science}.
        O'Reilly.
\end{itemize}

\end{document}
